\subsection{Risk management framework and governance}

These entries define risk management as a governed process and list the risks it must identify and the ways of measuring and changing them. Measurement methods are in \hyperref[meth:var-es-estimation-methods]{VaR and expected shortfall} and \hyperref[mod:factor-risk-model]{factor risk models}; the case numbers come from \nolinkurl{trading_notes.numerical.investment_process_planning.investment_planning_checks}.

\begin{qfdefinition}[Risk management framework: definition, process, risk governance, risk tolerance and risk budgeting; centralised and enterprise risk management; limits and capital allocation]\label{def:risk-management-framework}
\qffield{Statement} ``Risk management is a process involving the identification of exposures to risk, the establishment of appropriate ranges for exposures (given a clear understanding of an entity's objectives and constraints), the continuous measurement of these exposures (either present or contemplated), and the execution of appropriate adjustments whenever exposure levels fall outside of target ranges'' \citep[chap.~9 (Chance, Grant and Marsland), p.~580]{managing-investment-portfolios-2007}. The process starts by defining \emph{risk tolerance}, ``the level of risk it is willing and able to bear'', then identifies and measures risks, adjusts exposures when they leave target ranges, executes the risk-modifying transactions, and loops back to measurement (p.~581, Exhibit~9-1). Risk management ``involves adjusting levels of risk to appropriate levels, not necessarily eliminating risk altogether'', and can mean increasing risk where the firm has an edge (pp.~582--583). Tolerance therefore drives the rest of the process: it sets the target ranges and limits, decides which risks are retained and which are hedged or transferred, and defines the corrective action when measured risk is too high or too low; taking too little risk can be as damaging as too much (pp.~581 and~625).

\emph{Risk governance} is ``the process of setting overall policies and standards in risk management''. It involves choices of structure, infrastructure, reporting and methodology, and is judged by ``transparency, accountability, effectiveness \ldots and efficiency''. A centralised structure sees offsetting and diversified exposures across units and puts responsibility near senior management; this firm-wide view is enterprise risk management (p.~583). Effective governance for an investment firm separates the trading function from risk measurement and pricing and keeps the back office independent of the front office (p.~584).

\emph{Risk budgeting} allocates the overall risk budget among investments ``so as to maximize expected overall risk-adjusted return'' (chap.~1, p.~13). The same tools allocate capital across units through notional position limits, VaR-based limits, maximum loss limits and internal capital requirements. Each has a weakness: notional limits miss correlation, VaR limits are only as good as the VaR model and do not add up simply, and loss limits can be breached by market gaps (chap.~9, pp.~634--635).

\qffield{Practice} Interpretation, the features of a framework: (1) governance (board-approved policy, roles, independence); (2) a stated risk tolerance, ideally in the same units as the limits; (3) identification of all material risks, financial and non-financial; (4) measurement and stress testing; (5) limits and a risk budget; (6) monitoring and reporting; (7) a defined response when limits are breached (hedge, reduce, or approve the excess); (8) periodic review. Damodaran contrasts \emph{risk hedging}, which sees risk as a danger and buys protection with derivatives and insurance, with \emph{risk management}, which sees ``a danger and an opportunity'' and seeks advantage from uncertainty \citep[chap.~9, pp.~303--304, Table~9.7]{strategic-risk-taking-2008}.

In \hyperref[ex:investment-planning-case]{the planning case}, the tolerance is a 5th-percentile real loss before tax of $-13$ percent. The risk budget is visible in the Euler contributions: equity carries 91.7 percent of the balanced mix's variance on 55 percent of its capital. Allocation ranges serve as limits, and the IPS fixes the response to a breach.

\qffield{Conditions} The chapter is written for enterprises and portfolios alike (p.~579), and its regulatory examples (Basel~II) are dated to 2007. Risk tolerance may be multidimensional, with different tolerances for different risks that ``do not readily permit averaging'' (p.~581, n.~2).

\qffield{Pitfalls} A risk function that reports to the traders it monitors. Limits in units (notional) unrelated to the tolerance (loss). Treating zero risk as the goal: taking too little risk can be as costly as too much (p.~625). Adding unit VaRs as if they summed to firm VaR.

\qffield{See also} \hyperref[meth:risk-identification-modification]{identifying and modifying risks}; \hyperref[meth:investment-policy-statement]{IPS} (risk objective); \hyperref[meth:equal-risk-contribution]{risk contributions}.

\end{qfdefinition}

\begin{qfmethod}[Identifying, measuring and modifying risk exposures: financial and non-financial risks; measures; avoidance, acceptance and self-insurance, transfer by insurance, and shifting with derivatives]\label{meth:risk-identification-modification}
\qffield{Statement} Exposures are classified so that each can be measured and managed: market risk (interest rates, exchange rates, equity and commodity prices), credit risk, liquidity risk, operational risk, model risk, settlement risk, regulatory risk, legal or contract risk, tax risk, accounting risk and sovereign or political risk. ``Financial risk refers to all risks derived from events in the external financial markets; nonfinancial risk refers to all other forms of risk'' \citep[chap.~9, pp.~585--586, Exhibit~9-3]{managing-investment-portfolios-2007}. Market risk is the exposure of actively traded instruments to interest rates, exchange rates, equity and commodity prices, measured by techniques such as VaR and stress testing; credit and liquidity risk have their own measures. Non-financial risks ``are intrinsically very difficult to measure''. Some, such as regulatory, tax, legal and accounting risk, may not be measurable ``in any precise mathematical way'' because there is no observable loss distribution, and may suit insurance better than measurement and hedging (chap.~9, sec.~5, pp.~596 and~623).

\emph{Modifying exposures.} Damodaran lists the ways to reduce an exposure \citep[chap.~10, pp.~331--335]{strategic-risk-taking-2008}:
\begin{itemize}
\item through \emph{investment choices}, such as diversifying across locations or countries;
\item through \emph{financing choices}, such as borrowing in the currency of the cash flows;
\item by \emph{insurance}, which ``does not eliminate risk'' but ``shifts the risk from the firm buying the insurance to the insurance firm''. It is most effective against firm-specific, large and event risks; small risks can be self-insured;
\item with \emph{derivatives} (forwards, futures, options, swaps), which cover most market-wide risks.
\end{itemize}
\emph{Choosing among methods.} Forwards give the most complete, customised hedge when the cash flow is known but carry credit risk and higher cost for small flows; futures are cheaper and remove credit risk but need daily margin and may not match the exposure; options protect against the downside while keeping the upside, at a premium that grows with the protection. Giddy's rule uses forwards for known currency flows and options for unknown ones. Self-insurance against event risk makes sense when the firm can pool risks itself more cheaply than an insurer. ``The objective, after all, is not complete protection against risk, but as much protection as makes sense, given the marginal benefits and costs'' (pp.~338--339).

Interpretation, grouping these by what happens to the risk: \emph{avoidance} (do not take it, or exit); \emph{acceptance} (retain it deliberately, self-insure or diversify); \emph{transfer} (pay a premium for a third party to bear a loss); and \emph{shifting} (change the distribution of outcomes with derivatives or hedges, for example capping a loss with a put).

\emph{Interaction.} Risks are not independent: market losses can trigger liquidity risk when positions must be sold, leverage turns price moves into funding and counterparty risk, and model or operational weaknesses show up in stress. LTCM, leveraged about 25 times, ``failed to appreciate the market moves that would occur when it attempted to liquidate positions'', and its investors lost more than 90 percent \citep[chap.~9, p.~623]{managing-investment-portfolios-2007}. Interpretation: identification should therefore ask which risks rise together in the scenarios the investor fears.

\qffield{Practice} Example, from \hyperref[ex:investment-planning-case]{the planning case}. While the concentrated stock (22.2 percent of the pool, standard deviation 35 percent, correlation 0.6 with equity) is still held, portfolio volatility is 12.19 percent. Assuming the stock is expected to earn equity's 5.5 percent real, the pre-tax 5th percentile is $-16.43$ percent against a limit of $-13$. To cut equity risk quickly the client can sell equity index futures on 10 percent of the portfolio. The short future earns minus the equity return over cash, so with its cash leg the position equals moving 10 points from equity to cash. Volatility falls to 10.89 percent and the expected real return by 0.50 points, 10 percent of the 5.0-point equity premium over cash.

The hedge does not make the portfolio compliant: its 5th percentile is $-14.80$, and about 22 percent notional would be needed. An index future shifts only the stock's market exposure (beta 1.31); its specific risk, 28 of its 35 points of volatility, stays and can be removed only by selling. Combining the two methods works: selling the first third at once with the 10 percent hedge gives $-12.52$, and after the second sale the unhedged portfolio gives $-12.73$. Insurance does not suit a continuous market risk. Avoidance, selling everything at once, removes the specific risk immediately but brings forward the 60{,}000 of gains tax that staging defers. Shifting is therefore used for the transition, then retired.

\qffield{Conditions} Hedging changes exposure, not necessarily value; Damodaran finds the evidence that hedging raises firm value ``hazy at best'' (p.~331), and it costs premiums, spreads, basis risk and margin liquidity. The overlay example ignores basis, roll, the cash needed for variation margin and execution costs, and assumes the stock's beta and correlation are stable; a stock-specific shock is not hedged at all.

\qffield{Pitfalls} Measuring only market risk and missing operational, legal or model risk. Hedging a risk the investor is paid to bear. Treating a derivative hedge as free of counterparty and liquidity risk. Insuring small, frequent losses that could be self-insured.

\qffield{See also} \hyperref[def:risk-management-framework]{risk management framework}; \hyperref[meth:var-es-estimation-methods]{VaR and expected shortfall}.

\end{qfmethod}
